\documentclass[11pt,a4paper]{article}

% ============================================================
%  QIMMA -- COURS : Suites numériques
%  2ème Bac Sciences Mathématiques (A et B)
%  Standard exhaustif : définitions formelles + démonstrations
%  Basé sur templates/qimma-template.tex (charte Qimma)
% ============================================================

\usepackage[french]{babel}
\usepackage[utf8]{inputenc}
\usepackage[T1]{fontenc}
\usepackage{lmodern}
\usepackage[a4paper,margin=2cm,top=2.3cm,bottom=2.6cm]{geometry}
\usepackage{amsmath,amssymb}
\usepackage{xcolor}
\usepackage{graphicx}
\usepackage{tikz}
\usetikzlibrary{shadows,calc}
\usepackage[most]{tcolorbox}
\usepackage{titlesec}
\usepackage{fancyhdr}
\usepackage{enumitem}
\usepackage{pifont}
\usepackage{array}
\usepackage{colortbl}
\usepackage{hyperref}

% ------------------- CHARTE QIMMA -------------------
\definecolor{qteal}{HTML}{0d9488}
\definecolor{qtealdark}{HTML}{134e4a}
\definecolor{qamber}{HTML}{fbbf24}
\definecolor{qambertext}{HTML}{92400e}
\definecolor{ink}{HTML}{1f2937}
\definecolor{inkgray}{HTML}{6b7280}
\definecolor{tealpale}{HTML}{e6f5f3}
\colorlet{colDef}{qteal}
\definecolor{colProp}{HTML}{0f766e}
\definecolor{colEx}{HTML}{b45309}
\definecolor{colRem}{HTML}{9d174d}
\color{ink}

\graphicspath{{../../../../assets/}}
\newcommand{\logofile}{../../../../assets/qimma-logo.pdf}

% ------------------- METADONNEES -------------------
\newcommand{\matiere}{Mathématiques}
\newcommand{\niveau}{2\up{ème} Bac -- Sciences Mathématiques}
\newcommand{\chapitretitre}{Suites numériques}

% ------------------- EN-TETE / PIED DE PAGE -------------------
\pagestyle{fancy}
\fancyhf{}
\renewcommand{\headrulewidth}{0pt}
\renewcommand{\footrulewidth}{0.5pt}
\fancyfoot[L]{\raisebox{-0.15cm}{\includegraphics[height=0.35cm]{\logofile}}~\small\sffamily\color{inkgray}Qimma}
\fancyfoot[C]{\small\sffamily\color{inkgray}\matiere\ -- \niveau}
\fancyfoot[R]{\small\sffamily\color{qtealdark}\thepage}

% ------------------- TITRES DE SECTION -------------------
\titleformat{\section}
  {\normalfont\sffamily\Large\bfseries\color{qtealdark}}
  {\tikz[baseline=-3.2pt]{\node[circle,fill=qteal,text=white,inner sep=0pt,minimum size=8mm]{\sffamily\bfseries\thesection};}}
  {10pt}{}
  [\vspace{-2mm}{\color{qamber}\titlerule[1.6pt]}\vspace{2mm}]
\titlespacing*{\section}{0pt}{16pt}{10pt}

\titleformat{\subsection}
  {\normalfont\sffamily\large\bfseries\color{qtealdark}}
  {\color{qteal}\ding{212}}{6pt}{}
\titlespacing*{\subsection}{0pt}{12pt}{6pt}

% ------------------- BOITES PEDAGOGIQUES -------------------
\newtcolorbox{fiche}[3][]{
  enhanced, breakable,
  colback=white, colframe=#2,
  boxrule=1pt, arc=2mm,
  top=7mm, left=4mm, right=4mm, bottom=3mm,
  drop shadow={black!25, shadow xshift=1mm, shadow yshift=-1mm},
  attach boxed title to top left={xshift=4mm, yshift=-3mm, yshifttext=-1mm},
  boxed title style={colback=#2, arc=1.5mm, boxrule=0pt, left=2mm, right=2mm},
  coltitle=white, fonttitle=\sffamily\bfseries,
  title={#3}, #1
}
\newenvironment{definition}[1][Définition]
  {\begin{fiche}{colDef}{\ding{110}~~#1}}{\end{fiche}}
\newenvironment{propriete}[1][Propriété]
  {\begin{fiche}{colProp}{\ding{72}~~#1}}{\end{fiche}}
\newenvironment{exemple}[1][Exemple]
  {\begin{fiche}{colEx}{\ding{217}~~#1}}{\end{fiche}}
\newenvironment{remarque}[1][Remarque]
  {\begin{fiche}{colRem}{\ding{43}~~#1}}{\end{fiche}}

% Démonstration (hors boîte, style discret)
\newenvironment{demo}
  {\par\smallskip\noindent{\sffamily\bfseries\color{qtealdark}Démonstration.}\ \itshape}
  {\ \normalfont\hfill$\blacksquare$\par\medskip}

\hypersetup{colorlinks=true, linkcolor=qtealdark, urlcolor=qtealdark}

% ============================================================
\begin{document}

% ------------------- BANNIERE DE TITRE -------------------
\begin{tcolorbox}[
  enhanced, boxrule=0pt, arc=4mm,
  interior style={left color=qtealdark, right color=qteal},
  top=8mm, bottom=8mm, left=10mm, right=32mm,
  drop shadow={black!35, shadow xshift=1.5mm, shadow yshift=-1.5mm},
  before skip=0pt, after skip=9mm,
  overlay={
    \node[anchor=north west] at ([xshift=6mm,yshift=-6mm]frame.north west)
      {\includegraphics[height=1.25cm]{\logofile}};
    \node[anchor=north east, fill=qamber, text=qambertext, rounded corners=2pt,
          inner sep=4pt, font=\sffamily\bfseries\small] at ([xshift=-6mm,yshift=-6mm]frame.north east) {COURS};
  }
]
\hspace{1.6cm}\centering
{\sffamily\bfseries\color{white}\fontsize{23}{27}\selectfont \chapitretitre}\\[3mm]
{\sffamily\color{white!90}\large \matiere\ \textbullet\ \niveau}
\end{tcolorbox}

% ------------------- OBJECTIFS -------------------
\begin{tcolorbox}[
  enhanced, colback=tealpale, colframe=qteal, boxrule=0.8pt, arc=2mm,
  left=5mm, right=5mm, top=3mm, bottom=3mm,
  title={\sffamily\bfseries\color{qtealdark}\ding{51}~~Objectifs du chapitre},
  coltitle=qtealdark, colbacktitle=tealpale, boxrule=0.8pt,
  attach title to upper={\par\vspace{1mm}}
]
\begin{itemize}[leftmargin=6mm, label=\ding{212}, itemsep=1mm]
  \item Définir rigoureusement une suite, sa monotonie, ses bornes, à l'aide de quantificateurs.
  \item Maîtriser complètement les suites arithmétiques, géométriques et arithmético-géométriques (terme général, somme, limite).
  \item Démontrer et utiliser le théorème de la convergence monotone, le théorème des gendarmes et le théorème des suites adjacentes.
  \item Étudier une suite définie par récurrence $u_{n+1} = f(u_n)$ : intervalle stable, monotonie, convergence, limite.
  \item Utiliser la définition formelle de la convergence (en $\varepsilon$) et distinguer convergence / divergence vers $\pm\infty$ / divergence sans limite.
  \item Résoudre tous les types d'exercices type examen national sur les suites, y compris les suites couplées et les sommes télescopiques.
\end{itemize}
\end{tcolorbox}

\vspace{4mm}

% ------------------- SECTION 1 -------------------
\section{Généralités : définitions formelles}

\begin{definition}[Suite numérique]
Une \textbf{suite numérique} est une fonction $u : \mathbb{N} \to \mathbb{R}$ (ou définie à partir d'un certain rang $n_0$). L'image de $n$ par $u$ se note $u_n$ (plutôt que $u(n)$) et s'appelle le \textbf{terme général} (ou terme d'indice $n$) de la suite, notée $(u_n)_{n \geq n_0}$ ou simplement $(u_n)$.
\end{definition}

\begin{propriete}[Deux modes de définition]
\begin{enumerate}[leftmargin=7mm, itemsep=0.5mm]
  \item \textbf{Définition explicite} : $u_n = f(n)$, formule directe permettant de calculer n'importe quel terme sans connaître les précédents.
  \item \textbf{Définition par récurrence} : donnée du premier terme $u_0$ (ou $u_{n_0}$) et d'une relation $u_{n+1} = f(u_n)$ permettant de calculer chaque terme à partir du précédent.
\end{enumerate}
\end{propriete}

\begin{propriete}[Principe de récurrence]
Soit $P(n)$ une propriété dépendant de l'entier $n$, et $n_0\in\mathbb{N}$. Si :
\begin{itemize}[leftmargin=6mm, itemsep=0.5mm]
  \item \textbf{Initialisation} : $P(n_0)$ est vraie ;
  \item \textbf{Hérédité} : pour tout $n\geq n_0$, $\bigl(P(n) \text{ vraie}\bigr) \implies \bigl(P(n+1) \text{ vraie}\bigr)$ ;
\end{itemize}
alors $P(n)$ est vraie pour \textbf{tout} $n\geq n_0$.
\end{propriete}

\begin{remarque}[L'outil central de ce chapitre]
Le raisonnement par récurrence est \emph{la} méthode pour établir une propriété portant sur \textbf{tous} les termes d'une suite : formule explicite, encadrement, intervalle stable, monotonie... Il revient dans la quasi-totalité des démonstrations et exercices de ce chapitre ; assure-toi de toujours rédiger clairement les deux étapes (initialisation, hérédité) et de conclure explicitement.
\end{remarque}

\begin{definition}[Monotonie -- définitions avec quantificateurs]
La suite $(u_n)$ est :
\begin{itemize}[leftmargin=6mm, itemsep=0.5mm]
  \item \textbf{croissante} si $(\forall n \in \mathbb{N})\ u_{n+1} \geq u_n$ ; \textbf{strictement croissante} si $(\forall n)\ u_{n+1} > u_n$ ;
  \item \textbf{décroissante} si $(\forall n)\ u_{n+1} \leq u_n$ ; \textbf{strictement décroissante} si $(\forall n)\ u_{n+1} < u_n$ ;
  \item \textbf{constante} si $(\forall n)\ u_{n+1} = u_n$ ; \textbf{monotone} si elle est croissante ou décroissante.
\end{itemize}
\end{definition}

\begin{definition}[Suites majorées, minorées, bornées]
$(u_n)$ est \textbf{majorée} s'il existe $M \in \mathbb{R}$ tel que $(\forall n)\ u_n \leq M$ ; \textbf{minorée} s'il existe $m \in \mathbb{R}$ tel que $(\forall n)\ u_n \geq m$ ; \textbf{bornée} si elle est à la fois majorée et minorée, c'est-à-dire s'il existe $K \geq 0$ tel que $(\forall n)\ |u_n| \leq K$.
\end{definition}

\begin{propriete}[Trois méthodes pour étudier le sens de variation]
\begin{enumerate}[leftmargin=7mm, itemsep=0.5mm]
  \item \textbf{Signe de $u_{n+1} - u_n$} : méthode universelle. Si $u_{n+1} - u_n \geq 0$ pour tout $n$, la suite est croissante ; si $\leq 0$, décroissante.
  \item \textbf{Comparaison du rapport $\dfrac{u_{n+1}}{u_n}$ à $1$} : valable \textbf{uniquement} si tous les termes sont strictement positifs (ou tous strictement négatifs, avec inégalité inversée).
  \item \textbf{Étude de la fonction associée} : si $u_n = f(n)$ avec $f$ monotone sur $[0\,;+\infty[$, la suite a le même sens de variation que $f$.
\end{enumerate}
\end{propriete}

\begin{exemple}[Les trois méthodes en pratique]
\textbf{a)} $u_n = n^2 - 5n$ : $u_{n+1} - u_n = (n+1)^2 - 5(n+1) - n^2 + 5n = 2n - 4$. Ce signe dépend de $n$ (négatif si $n < 2$, positif si $n \geq 2$) : $(u_n)$ n'est \textbf{pas monotone} sur $\mathbb{N}$ entier, mais décroissante puis croissante (minimum en $n=2$).

\smallskip
\textbf{b)} $u_n = \dfrac{3^n}{n!}$ (tous strictement positifs) : $\dfrac{u_{n+1}}{u_n} = \dfrac{3}{n+1}$. Pour $n \geq 3$, ce rapport est $\leq 1$ : $(u_n)$ est décroissante à partir de $n = 3$.

\smallskip
\textbf{c)} $u_n = \sqrt{n+1}$ : $f(x) = \sqrt{x+1}$ est croissante sur $[0\,;+\infty[$, donc $(u_n)$ est croissante.
\end{exemple}

% ------------------- SECTION 2 -------------------
\section{Suites arithmétiques et géométriques}

\subsection{Suites arithmétiques}

\begin{definition}
$(u_n)$ est \textbf{arithmétique} de raison $r \in \mathbb{R}$ si $(\forall n)\ u_{n+1} = u_n + r$.
\end{definition}

\begin{propriete}[Terme général et somme]
\[
  u_n = u_0 + n\,r \qquad\text{(plus généralement } u_n = u_p + (n-p)\,r\text{)}.
\]
\end{propriete}

\begin{demo}
Par récurrence sur $n$. \textbf{Initialisation} : $u_0 = u_0 + 0 \times r$, vrai. \textbf{Hérédité} : supposons $u_n = u_0 + nr$. Alors $u_{n+1} = u_n + r = u_0 + nr + r = u_0 + (n+1)r$. La propriété est vraie au rang $n+1$, donc pour tout $n \in \mathbb{N}$.
\end{demo}

\begin{propriete}[Somme de termes consécutifs]
\[
  S = u_p + u_{p+1} + \cdots + u_q = (q - p + 1) \times \frac{u_p + u_q}{2}
\]
(nombre de termes $\times$ moyenne du premier et du dernier). Cas particulier fondamental : $1 + 2 + \cdots + n = \dfrac{n(n+1)}{2}$.
\end{propriete}

\begin{demo}
Notons $S = u_p + u_{p+1} + \cdots + u_q$. En écrivant la somme à l'envers, $S = u_q + u_{q-1} + \cdots + u_p$, puis en additionnant terme à terme les deux écritures :
\[
  2S = \sum_{k=p}^{q} (u_k + u_{q+p-k}) = (q - p + 1)(u_p + u_q),
\]
car chaque paire $u_k + u_{q+p-k}$ vaut $u_p + u_q$ (les termes étant en progression arithmétique). D'où $S = (q-p+1)\dfrac{u_p+u_q}{2}$.
\end{demo}

\subsection{Suites géométriques}

\begin{definition}
$(u_n)$ est \textbf{géométrique} de raison $q \in \mathbb{R}$ si $(\forall n)\ u_{n+1} = q\,u_n$.
\end{definition}

\begin{propriete}[Terme général et somme]
\[
  u_n = u_0 \cdot q^n \qquad\text{(plus généralement } u_n = u_p \cdot q^{\,n-p}\text{)}.
\]
Pour $q \neq 1$ : \[S = u_p + u_{p+1} + \cdots + u_q = u_p \times \frac{1 - q^{\,q-p+1}}{1 - q}.\]
Si $q = 1$ : $S = (q - p + 1)\,u_p$ (tous les termes sont égaux).
\end{propriete}

\begin{demo}
Notons $S = u_p(1 + q + \cdots + q^{q-p})$. En posant $n = q - p$ et $T = 1 + q + \cdots + q^n$ :
\[
  T - qT = (1 + q + \cdots + q^n) - (q + q^2 + \cdots + q^{n+1}) = 1 - q^{n+1},
\]
donc $T(1-q) = 1 - q^{n+1}$, d'où $T = \dfrac{1-q^{n+1}}{1-q}$ si $q \neq 1$, et $S = u_p T$.
\end{demo}

\begin{exemple}[Reconnaître et calculer]
Soit $u_0 = 3$ et $u_{n+1} = -2u_n$. Calculons $S = u_0 + u_1 + \cdots + u_5$.

\smallskip
$(u_n)$ est géométrique de raison $q = -2$, premier terme $u_0 = 3$ :
\[
  S = 3 \times \frac{1 - (-2)^6}{1 - (-2)} = 3 \times \frac{1 - 64}{3} = 1 - 64 = -63.
\]
\end{exemple}

\subsection{Suites arithmético-géométriques}

\begin{propriete}[Méthode du point fixe -- $u_{n+1} = a\,u_n + b$ avec $a \neq 1$]
\begin{enumerate}[leftmargin=7mm, itemsep=0.5mm]
  \item \textbf{Point fixe} : résoudre $\ell = a\ell + b$, soit $\ell = \dfrac{b}{1-a}$.
  \item \textbf{Suite auxiliaire} : poser $v_n = u_n - \ell$. Alors $v_{n+1} = u_{n+1} - \ell = au_n + b - \ell = a(u_n - \ell) = a\,v_n$ (car $\ell = a\ell+b$) : $(v_n)$ est géométrique de raison $a$.
  \item \textbf{Terme général} : $v_n = v_0\,a^n$, d'où $u_n = \ell + (u_0 - \ell)\,a^n$.
  \item \textbf{Limite} : si $|a| < 1$, $a^n \to 0$ donc $u_n \to \ell$.
\end{enumerate}
\end{propriete}

\begin{exemple}[Application complète]
Soit $u_0 = 5$, $u_{n+1} = \dfrac{1}{2}u_n + 3$. Exprimer $u_n$ en fonction de $n$ et calculer $\lim u_n$.

\smallskip
Point fixe : $\ell = \frac{1}{2}\ell + 3 \Rightarrow \ell = 6$. Suite auxiliaire $v_n = u_n - 6$ : géométrique de raison $\frac{1}{2}$, $v_0 = -1$. Donc $v_n = -\left(\frac{1}{2}\right)^n$, et $u_n = 6 - \left(\frac{1}{2}\right)^n$. Comme $\left|\frac{1}{2}\right| < 1$ : $\lim u_n = 6$.
\end{exemple}

% ------------------- SECTION 3 -------------------
\section{Convergence : définition formelle et propriétés}

\begin{definition}[Limite finie d'une suite -- définition en $\varepsilon$]
$(u_n)$ \textbf{converge} vers $\ell \in \mathbb{R}$, noté $\lim\limits_{n \to +\infty} u_n = \ell$, lorsque :
\[
  (\forall \varepsilon > 0)\ (\exists N \in \mathbb{N})\ (\forall n \geq N)\quad |u_n - \ell| < \varepsilon.
\]
Une suite non convergente est dite \textbf{divergente} (elle tend vers $+\infty$, vers $-\infty$, ou n'a pas de limite).
\end{definition}

\begin{propriete}[Unicité de la limite]
Si $(u_n)$ converge, sa limite est \textbf{unique}.
\end{propriete}

\begin{demo}
Identique à la preuve pour les fonctions : si $u_n \to \ell$ et $u_n \to \ell'$ avec $\ell \neq \ell'$, on pose $\varepsilon = \frac{|\ell-\ell'|}{2}$ ; pour $n$ assez grand, $|u_n - \ell| < \varepsilon$ et $|u_n - \ell'| < \varepsilon$, donc $|\ell - \ell'| \leq |\ell - u_n| + |u_n - \ell'| < 2\varepsilon = |\ell - \ell'|$, absurde.
\end{demo}

\begin{propriete}[Limites infinies]
\[
  \lim_{n \to +\infty} u_n = +\infty \iff (\forall A > 0)(\exists N)(\forall n \geq N)\ u_n > A,
\]
et de façon analogue pour $-\infty$.
\end{propriete}

\begin{propriete}[Opérations sur les limites]
Les mêmes règles que pour les fonctions en $+\infty$ s'appliquent : somme, produit, quotient des limites (avec les mêmes conventions et les mêmes \textbf{formes indéterminées} $\infty-\infty$, $0\times\infty$, $\frac{\infty}{\infty}$, $\frac{0}{0}$). Toute suite convergente est \textbf{bornée} (réciproque fausse : $u_n = (-1)^n$ est bornée mais divergente).
\end{propriete}

\begin{exemple}[Limites de référence pour $q^n$]
\begin{itemize}[leftmargin=8mm, itemsep=0.5mm]
  \item si $-1 < q < 1$ : $q^n \to 0$ ;
  \item si $q > 1$ : $q^n \to +\infty$ ;
  \item si $q = 1$ : $q^n = 1$ pour tout $n$ ;
  \item si $q \leq -1$ : $(q^n)$ diverge sans limite (oscillations d'amplitude croissante ou constante).
\end{itemize}
\end{exemple}

% ------------------- SECTION 4 -------------------
\section{Théorèmes de convergence}

\begin{propriete}[Théorème de la convergence monotone -- admis]
\begin{itemize}[leftmargin=6mm, itemsep=0.5mm]
  \item Toute suite \textbf{croissante et majorée} converge (vers un réel $\ell \leq$ tout majorant).
  \item Toute suite \textbf{décroissante et minorée} converge.
  \item Toute suite \textbf{croissante et non majorée} tend vers $+\infty$.
  \item Toute suite \textbf{décroissante et non minorée} tend vers $-\infty$.
\end{itemize}
\emph{Ce théorème garantit l'\textbf{existence} de la limite, mais ne donne pas sa valeur.}
\end{propriete}

\begin{propriete}[Théorème des gendarmes pour les suites]
Si $v_n \leq u_n \leq w_n$ à partir d'un certain rang, et si $\lim v_n = \lim w_n = \ell$, alors $\lim u_n = \ell$.

\smallskip
\textbf{Variante utile} : si $|u_n - \ell| \leq v_n$ et $v_n \to 0$, alors $u_n \to \ell$.
\end{propriete}

\begin{propriete}[Théorème de comparaison]
Si $u_n \geq v_n$ à partir d'un certain rang et $v_n \to +\infty$, alors $u_n \to +\infty$. De même, si $u_n \leq v_n$ et $v_n \to -\infty$, alors $u_n \to -\infty$.
\end{propriete}

\begin{exemple}[Convergence monotone -- schéma complet]
Soit $u_0 = 0$ et $u_{n+1} = \sqrt{2 + u_n}$. Montrons que $(u_n)$ converge et calculons sa limite.

\smallskip
\textbf{Étape 1 (intervalle stable)} : montrons par récurrence que $0 \leq u_n \leq 2$ pour tout $n$. Initialisation : $u_0 = 0 \in [0\,;2]$. Hérédité : si $0 \leq u_n \leq 2$, alors $2 \leq u_n + 2 \leq 4$, donc $\sqrt{2} \leq u_{n+1} \leq 2$, en particulier $0 \leq u_{n+1} \leq 2$.

\smallskip
\textbf{Étape 2 (monotonie)} : $u_{n+1} - u_n = \sqrt{2+u_n} - u_n$. On étudie $g(x) = \sqrt{2+x} - x$ sur $[0\,;2]$ : $g'(x) = \frac{1}{2\sqrt{2+x}} - 1 < 0$ (car $2\sqrt{2+x} \geq 2\sqrt{2} > 1$), donc $g$ est décroissante, et $g(2) = \sqrt{4} - 2 = 0$. Ainsi $g(x) \geq g(2) = 0$ pour $x \leq 2$ : $u_{n+1} - u_n = g(u_n) \geq 0$, la suite est \textbf{croissante}.

\smallskip
\textbf{Étape 3 (convergence)} : $(u_n)$ est croissante et majorée par $2$ : elle converge vers un réel $\ell \in [0\,;2]$.

\smallskip
\textbf{Étape 4 (calcul de $\ell$)} : $f(x) = \sqrt{2+x}$ est continue sur $[0\,;2]$, donc $\ell = \sqrt{2+\ell}$, soit $\ell^2 = 2 + \ell$ (car $\ell \geq 0$), c'est-à-dire $\ell^2 - \ell - 2 = 0$, $(\ell-2)(\ell+1)=0$. Comme $\ell \geq 0$ : $\boxed{\ell = 2}$.
\end{exemple}

\begin{remarque}[Rédaction attendue -- les 4 étapes obligatoires]
Pour une suite récurrente $u_{n+1} = f(u_n)$, la rédaction complète exige \textbf{toujours} : 1) un intervalle stable par $f$ contenant tous les $u_n$ (par récurrence) ; 2) le sens de variation ; 3) l'application du théorème de convergence monotone ; 4) la résolution de $f(\ell) = \ell$ \emph{après} avoir justifié l'existence de $\ell$. Sauter une étape fait perdre des points, même si le résultat final est correct.
\end{remarque}

% ------------------- SECTION 5 -------------------
\section{Suites adjacentes}

\begin{definition}[Suites adjacentes]
Deux suites $(u_n)$ et $(v_n)$ sont \textbf{adjacentes} lorsque :
\begin{enumerate}[leftmargin=7mm, itemsep=0.5mm]
  \item l'une est croissante et l'autre décroissante (disons $(u_n)$ croissante et $(v_n)$ décroissante) ;
  \item $\lim\limits_{n\to+\infty}(v_n-u_n)=0$.
\end{enumerate}
\end{definition}

\begin{propriete}[Théorème des suites adjacentes]
Si $(u_n)$ et $(v_n)$ sont adjacentes ($(u_n)$ croissante, $(v_n)$ décroissante), alors :
\begin{itemize}[leftmargin=6mm, itemsep=0.5mm]
  \item pour tout $n$ : $u_n\leq v_n$ (et même $u_0\leq u_n\leq v_n\leq v_0$) ;
  \item $(u_n)$ et $(v_n)$ \textbf{convergent}, et vers la \textbf{même limite} $\ell$, avec $u_n\leq\ell\leq v_n$ pour tout $n$.
\end{itemize}
\end{propriete}

\begin{demo}
Posons $w_n=v_n-u_n$. Alors $w_{n+1}-w_n=(v_{n+1}-v_n)-(u_{n+1}-u_n)\leq0$ (car $v_{n+1}-v_n\leq0$ et $u_{n+1}-u_n\geq0$) : $(w_n)$ est \textbf{décroissante}. Comme de plus $w_n\to0$, une suite décroissante converge vers sa borne inférieure, donc $w_n\geq\lim w_n=0$ pour tout $n$ : ainsi $v_n\geq u_n$ pour tout $n$.

\smallskip
Puisque $u_n\leq v_n\leq v_0$ (car $(v_n)$ décroissante), $(u_n)$ est croissante et \textbf{majorée} par $v_0$ : elle converge vers un réel $\ell_1$ (théorème de la convergence monotone). De même, $(v_n)$ est décroissante et minorée par $u_0$ : elle converge vers $\ell_2$.

\smallskip
Or $v_n-u_n\to\ell_2-\ell_1$ (limite d'une différence de suites convergentes) et $v_n-u_n\to0$ par hypothèse : par unicité de la limite, $\ell_2-\ell_1=0$, soit $\ell_1=\ell_2=\ell$.

\smallskip
Enfin, $(u_n)$ étant croissante vers $\ell$ : $u_n\leq\ell$ pour tout $n$ ; $(v_n)$ étant décroissante vers $\ell$ : $\ell\leq v_n$ pour tout $n$. D'où $u_n\leq\ell\leq v_n$.
\end{demo}

\begin{remarque}[Pourquoi ce théorème est puissant]
Il permet de prouver la convergence de \textbf{deux} suites \textbf{simultanément}, sans connaître à l'avance leur limite commune, et fournit à chaque rang un \textbf{encadrement} $u_n\leq\ell\leq v_n$ -- exactement le principe de la \textbf{dichotomie} vue au chapitre « Limites et continuité » : les suites $(a_n)$ et $(b_n)$ construites par dichotomie pour approcher une solution de $f(x)=0$ sont en fait deux suites \textbf{adjacentes} (l'une croissante, l'autre décroissante, d'écart $\frac{b-a}{2^n}\to0$).
\end{remarque}

\begin{exemple}[Exemple classique de suites adjacentes]
Soient $u_n=\displaystyle\sum_{k=1}^n\frac{1}{k^2}$ et $v_n=u_n+\dfrac1n$, pour $n\geq1$. Montrons qu'elles sont adjacentes.

\smallskip
\emph{Monotonie de $(u_n)$} : $u_{n+1}-u_n=\dfrac{1}{(n+1)^2}>0$ : $(u_n)$ est strictement croissante.

\smallskip
\emph{Monotonie de $(v_n)$} :
\[
  v_{n+1}-v_n = \underbrace{u_{n+1}-u_n}_{\frac{1}{(n+1)^2}} + \frac{1}{n+1}-\frac1n
  = \frac1{(n+1)^2} + \frac{n-(n+1)}{n(n+1)}
  = \frac1{(n+1)^2} - \frac1{n(n+1)}.
\]
En réduisant au même dénominateur : $v_{n+1}-v_n=\dfrac{n-(n+1)}{n(n+1)^2}=\dfrac{-1}{n(n+1)^2}<0$ : $(v_n)$ est strictement décroissante.

\smallskip
\emph{Écart} : $v_n-u_n=\dfrac1n\xrightarrow[n\to+\infty]{}0$.

\smallskip
\textbf{Conclusion} : $(u_n)$ et $(v_n)$ sont adjacentes, donc convergent vers une même limite $\ell$, avec $u_n\leq\ell\leq v_n$ pour tout $n$ (on démontre, hors-programme, que $\ell=\frac{\pi^2}{6}$).
\end{exemple}

% ------------------- SECTION 6 -------------------
\section{Convergence par l'inégalité des accroissements finis}

\begin{propriete}[Méthode alternative de convergence]
Soit $(u_n)$ définie par $u_{n+1} = f(u_n)$, $\ell$ un point fixe de $f$ ($f(\ell)=\ell$), et $I$ un intervalle stable par $f$ contenant tous les $u_n$ et $\ell$. Si $f$ est dérivable sur $I$ avec $|f'(x)| \leq k < 1$ pour tout $x \in I$, alors :
\[
  |u_n - \ell| \leq k^n\,|u_0 - \ell| \xrightarrow[n \to +\infty]{} 0,
  \qquad\text{donc } u_n \to \ell.
\]
\end{propriete}

\begin{demo}
Par l'inégalité des accroissements finis entre $u_n$ et $\ell$ (tous deux dans $I$) :
\[
  |u_{n+1} - \ell| = |f(u_n) - f(\ell)| \leq k\,|u_n - \ell|.
\]
Par récurrence immédiate : $|u_n - \ell| \leq k^n |u_0 - \ell|$. Comme $0 \leq k < 1$, $k^n \to 0$, donc par encadrement $|u_n - \ell| \to 0$, c'est-à-dire $u_n \to \ell$.
\end{demo}

\begin{remarque}[Avantage de cette méthode]
Elle \textbf{évite} d'étudier la monotonie de $(u_n)$ (parfois délicate) et donne \textbf{en prime} une majoration de la vitesse de convergence, utile pour les questions du type « déterminer $n$ tel que $|u_n - \ell| < 10^{-p}$ ».
\end{remarque}

% ------------------- SECTION 7 -------------------
\section{Sommes télescopiques et suites couplées}

\begin{propriete}[Somme télescopique]
Si $u_n = a_{n+1} - a_n$ pour une suite auxiliaire $(a_n)$, alors :
\[
  \sum_{k=0}^{n} u_k = \sum_{k=0}^{n} (a_{k+1} - a_k) = a_{n+1} - a_0
\]
(tous les termes intermédiaires s'annulent deux à deux).
\end{propriete}

\begin{exemple}[Somme télescopique classique]
Calculons $S_n = \displaystyle\sum_{k=1}^{n} \frac{1}{k(k+1)}$.

\smallskip
Décomposition : $\dfrac{1}{k(k+1)} = \dfrac{1}{k} - \dfrac{1}{k+1}$ (à vérifier par réduction au même dénominateur). Donc
\[
  S_n = \sum_{k=1}^{n} \left(\frac{1}{k} - \frac{1}{k+1}\right) = \frac{1}{1} - \frac{1}{n+1} = 1 - \frac{1}{n+1} \xrightarrow[n \to +\infty]{} 1.
\]
\end{exemple}

\begin{remarque}[Suites couplées]
Certains exercices définissent deux suites liées, par exemple $u_{n+1} = \frac{u_n+v_n}{2}$ et $v_{n+1} = \sqrt{u_n v_n}$ (moyenne arithmético-géométrique), ou $u_{n+1} = 2u_n - v_n$, $v_{n+1} = u_n$. La méthode consiste souvent à étudier une combinaison $w_n = u_n - v_n$ ou $w_n = u_n + v_n$ pour se ramener à une suite arithmétique ou géométrique connue.
\end{remarque}

% ------------------- SECTION 8 -------------------
\section{Exercices corrigés -- niveau examen national SM}

\begin{exemple}[Exercice 1 -- arithmético-géométrique complète]
Soit $u_0 = 1$ et $u_{n+1} = 3u_n - 4$. Montrer que $v_n = u_n - 2$ est géométrique, exprimer $u_n$, étudier la monotonie et la limite.

\smallskip
\textbf{Corrigé.} $v_{n+1} = u_{n+1} - 2 = 3u_n - 6 = 3(u_n - 2) = 3v_n$ : $(v_n)$ géométrique de raison $3$, $v_0 = -1$, donc $v_n = -3^n$, d'où $u_n = 2 - 3^n$.

\smallskip
Monotonie : $u_{n+1} - u_n = -3^{n+1} + 3^n = 3^n(1 - 3) = -2\times 3^n < 0$ : $(u_n)$ est strictement décroissante.

\smallskip
Limite : $3^n \to +\infty$, donc $u_n = 2 - 3^n \to -\infty$.
\end{exemple}

\begin{exemple}[Exercice 2 -- suite définie par une somme]
Soit $S_n = \displaystyle\sum_{k=0}^{n} \frac{1}{2^k}$. Calculer $S_n$ en fonction de $n$, puis $\lim S_n$.

\smallskip
\textbf{Corrigé.} $S_n$ est la somme des $n+1$ premiers termes d'une suite géométrique de premier terme $1$ et de raison $\frac{1}{2}$ :
\[
  S_n = \frac{1 - \left(\frac{1}{2}\right)^{n+1}}{1 - \frac{1}{2}} = 2\left(1 - \left(\frac{1}{2}\right)^{n+1}\right) = 2 - \left(\frac{1}{2}\right)^n.
\]
Comme $\left(\frac{1}{2}\right)^n \to 0$ : $\lim S_n = 2$.
\end{exemple}

\begin{exemple}[Exercice 3 -- convergence monotone avec fonction homographique]
Soit $u_0 = 2$ et $u_{n+1} = \dfrac{2u_n + 1}{u_n + 2}$. Montrer que $(u_n)$ est bien définie, à valeurs dans $[1\,;2]$, décroissante, et calculer sa limite.

\smallskip
\textbf{Corrigé.} Posons $f(x) = \dfrac{2x+1}{x+2}$ sur $[1\,;2]$ ($x + 2 \neq 0$, la suite est bien définie tant que $u_n \neq -2$, ce qui sera garanti par $u_n \in [1\,;2]$).

\smallskip
\emph{Intervalle stable} : $f(1) = 1$, $f(2) = \frac{5}{4}$ ; $f'(x) = \frac{2(x+2)-(2x+1)}{(x+2)^2} = \frac{3}{(x+2)^2} > 0$ : $f$ est croissante sur $[1\,;2]$, donc $f([1\,;2]) = [f(1)\,;f(2)] = [1\,;1{,}25] \subset [1\,;2]$. Par récurrence, $u_n \in [1\,;2]$ pour tout $n$.

\smallskip
\emph{Monotonie} : $u_{n+1} - u_n = f(u_n) - u_n = \dfrac{2u_n+1 - u_n(u_n+2)}{u_n+2} = \dfrac{1 - u_n^2}{u_n+2}$. Pour $u_n \in [1\,;2]$ : $1 - u_n^2 \leq 0$ et $u_n + 2 > 0$, donc $u_{n+1} - u_n \leq 0$ : $(u_n)$ est décroissante.

\smallskip
\emph{Convergence} : décroissante et minorée par $1$ : converge vers $\ell \in [1\,;2]$. $f$ continue en $\ell$ : $\ell = \dfrac{2\ell+1}{\ell+2}$, soit $\ell(\ell+2) = 2\ell+1$, $\ell^2 = 1$, donc $\ell = \pm 1$. Comme $\ell \geq 1$ : $\boxed{\ell = 1}$.
\end{exemple}

\begin{exemple}[Exercice 4 -- deux suites couplées, adjacentes]
Soit $(u_n)$ et $(v_n)$ définies par $u_0=1$, $v_0=7$, et pour tout $n$ : $u_{n+1} = \dfrac{u_n+v_n}{2}$, $v_{n+1} = \dfrac{u_n+3v_n}{4}$. Montrer que $(u_n)$ et $(v_n)$ sont adjacentes, puis déterminer leur limite commune.

\smallskip
\textbf{Corrigé.} \emph{Écart} : posons $w_n=v_n-u_n$.
\[
  w_{n+1} = v_{n+1} - u_{n+1} = \frac{u_n+3v_n}{4} - \frac{u_n+v_n}{2} = \frac{u_n+3v_n - 2u_n - 2v_n}{4} = \frac{v_n - u_n}{4} = \frac14w_n.
\]
$(w_n)$ est géométrique de raison $\frac14$, $w_0=7-1=6$, donc $w_n=6\left(\frac14\right)^n\to0$ ; de plus $w_n>0$ pour tout $n$ (produit de termes positifs), donc $v_n>u_n$ pour tout $n$.

\smallskip
\emph{Monotonie de $(u_n)$} : $u_{n+1}-u_n=\dfrac{u_n+v_n}{2}-u_n=\dfrac{v_n-u_n}{2}=\dfrac{w_n}{2}>0$ : $(u_n)$ est \textbf{croissante}.

\smallskip
\emph{Monotonie de $(v_n)$} : $v_{n+1}-v_n=\dfrac{u_n+3v_n}{4}-v_n=\dfrac{u_n-v_n}{4}=\dfrac{-w_n}{4}<0$ : $(v_n)$ est \textbf{décroissante}.

\smallskip
Comme $(u_n)$ croissante, $(v_n)$ décroissante et $v_n-u_n=w_n\to0$ : $(u_n)$ et $(v_n)$ sont \textbf{adjacentes}, donc convergent vers une \textbf{même limite} $\ell$.

\smallskip
\emph{Détermination de $\ell$} : cherchons une combinaison $au_n+bv_n$ constante. On veut $a\cdot\dfrac{u_n+v_n}{2}+b\cdot\dfrac{u_n+3v_n}{4}=au_n+bv_n$, soit $\dfrac{2a+b}{4}=a$ et $\dfrac{2a+3b}{4}=b$, ce qui donne $b=2a$ dans les deux cas. Avec $a=1$, $b=2$ : $u_n+2v_n$ est \textbf{constante}.
\[
  u_n+2v_n = u_0+2v_0 = 1+14 = 15 \qquad \text{pour tout } n.
\]
Par passage à la limite ($u_n\to\ell$, $v_n\to\ell$) : $\ell+2\ell=15$, soit $\boxed{\ell=5}$.
\end{exemple}

\begin{exemple}[Exercice 5 -- majoration de la vitesse de convergence]
Reprenons $u_0=0$, $u_{n+1}=\sqrt{2+u_n}$ (limite $\ell=2$ trouvée plus haut). Montrer que $|u_{n+1}-2| \leq \frac{1}{2}|u_n - 2|$, puis déterminer un entier $n$ tel que $|u_n - 2| < 10^{-3}$.

\smallskip
\textbf{Corrigé.} $f(x) = \sqrt{2+x}$, $f'(x) = \dfrac{1}{2\sqrt{2+x}}$. Sur $[0\,;2]$, $2+x \in [2\,;4]$, donc $\sqrt{2+x} \geq \sqrt{2}$, d'où $f'(x) \leq \dfrac{1}{2\sqrt{2}} < \dfrac{1}{2}$.

\smallskip
Par l'inégalité des accroissements finis : $|u_{n+1}-2| = |f(u_n)-f(2)| \leq \frac{1}{2}|u_n-2|$, donc $|u_n - 2| \leq \left(\frac{1}{2}\right)^n |u_0-2| = 2\left(\frac{1}{2}\right)^n$.

\smallskip
On veut $2\left(\frac{1}{2}\right)^n < 10^{-3}$, soit $\left(\frac{1}{2}\right)^{n-1} < 10^{-3}$, soit $n - 1 > \dfrac{3\ln 10}{\ln 2} \approx 9{,}97$, donc $n \geq 11$ convient.
\end{exemple}

\begin{exemple}[Exercice 6 -- suite non monotone, encadrement direct]
Soit $u_n = \dfrac{(-1)^n}{n} + 3$ pour $n \geq 1$. Montrer que $(u_n)$ converge vers $3$ sans étudier sa monotonie (elle n'est pas monotone).

\smallskip
\textbf{Corrigé.} Pour tout $n \geq 1$ : $|u_n - 3| = \left|\dfrac{(-1)^n}{n}\right| = \dfrac{1}{n} \to 0$. Par la variante du théorème des gendarmes ($|u_n - \ell| \leq v_n \to 0 \Rightarrow u_n \to \ell$) : $\lim u_n = 3$.
\end{exemple}

% ------------------- SECTION 9 -------------------
\section{Méthodes et pièges : la boîte à outils de l'examen}

\subsection{Ce qu'on te demande, ce que tu fais}

\begin{center}
\renewcommand{\arraystretch}{1.45}
\sffamily\small
\begin{tabular}{>{\raggedright\arraybackslash}p{7.2cm} >{\raggedright\arraybackslash}p{8cm}}
\rowcolor{qtealdark} \color{white}\bfseries Ce qu'on te demande & \color{white}\bfseries Ton réflexe \\
\rowcolor{tealpale} Montrer qu'une suite est arithmétique/géométrique & Calculer $u_{n+1}-u_n$ (constant $\Rightarrow$ arithm.) ou $\frac{u_{n+1}}{u_n}$ (constant $\Rightarrow$ géom.). \\
Suite arithmético-géométrique $u_{n+1}=au_n+b$ & Point fixe $\ell=\frac{b}{1-a}$, puis $v_n=u_n-\ell$ géométrique de raison $a$. \\
\rowcolor{tealpale} Étudier la monotonie & Signe de $u_{n+1}-u_n$ (universel) ; rapport si termes positifs ; fonction associée si $u_n=f(n)$. \\
Prouver la convergence d'une récurrente & 4 étapes : intervalle stable, monotonie, convergence monotone, résoudre $f(\ell)=\ell$. \\
\rowcolor{tealpale} Convergence sans étudier la monotonie & Inégalité des accroissements finis : $|f'|\leq k<1$ sur un intervalle stable. \\
Trouver $n$ tel que $|u_n-\ell|<10^{-p}$ & Utiliser la majoration $|u_n-\ell|\leq k^n|u_0-\ell|$, résoudre l'inéquation avec $\ln$. \\
\rowcolor{tealpale} Prouver la convergence de 2 suites liées sans connaître $\ell$ & Montrer qu'elles sont \textbf{adjacentes} (monotonies opposées + écart $\to0$). \\
Trouver la limite commune de 2 suites adjacentes & Chercher une combinaison linéaire $au_n+bv_n$ \textbf{constante}, puis passer à la limite. \\
\rowcolor{tealpale} Calculer une somme $\sum u_k$ & Reconnaître arithmétique/géométrique, ou chercher une écriture télescopique $a_{k+1}-a_k$. \\
Suites couplées $(u_n),(v_n)$ & Étudier une combinaison $w_n = u_n \pm v_n$ ou $u_n - v_n$ pour retrouver une suite connue (souvent adjacentes). \\
\rowcolor{tealpale} Suite non monotone qui converge quand même & Encadrer $|u_n - \ell|$ par une suite qui tend vers $0$ (gendarmes). \\
\end{tabular}
\end{center}

\subsection{Les pièges qui coûtent des points}

\begin{remarque}[Erreurs relevées chaque année par les correcteurs]
\begin{itemize}[leftmargin=6mm, itemsep=0.5mm]
  \item Conclure $\lim u_n = \ell$ en résolvant $f(\ell)=\ell$ \textbf{sans avoir prouvé la convergence} au préalable (monotonie+borne, ou AF).
  \item Utiliser le \textbf{rapport} $\frac{u_{n+1}}{u_n}$ pour la monotonie alors que les termes ne sont pas tous positifs.
  \item Oublier de démontrer par \textbf{récurrence} que la suite reste dans l'intervalle stable (l'affirmer sans preuve).
  \item Confondre « $(u_n)$ bornée » et « $(u_n)$ convergente » : une suite bornée peut diverger (ex. $(-1)^n$).
  \item Appliquer la formule de somme géométrique avec $q=1$ (division par $0$) : traiter ce cas à part.
  \item Oublier le facteur $(q-p+1)$ (nombre de termes) dans la somme arithmétique ou géométrique entre les indices $p$ et $q$.
  \item Pour l'inégalité des accroissements finis : oublier de vérifier que $k<1$ \textbf{strictement} (sinon $k^n$ ne tend pas vers $0$).
  \item Pour les suites adjacentes : oublier de préciser \textbf{laquelle} est croissante et laquelle est décroissante, ou conclure « adjacentes » sans avoir vérifié \textbf{les deux} conditions (monotonies opposées \emph{et} écart tendant vers $0$).
  \item Croire que le théorème des suites adjacentes donne la \textbf{valeur} de la limite commune : il ne donne que son \textbf{existence} et l'encadrement $u_n\leq\ell\leq v_n$ ; il faut un argument supplémentaire (combinaison linéaire constante, etc.) pour calculer $\ell$.
\end{itemize}
\end{remarque}

% ------------------- RESUME -------------------
\vspace{4mm}
\begin{tcolorbox}[
  enhanced, boxrule=0pt, arc=2mm,
  interior style={left color=qtealdark, right color=qteal},
  left=6mm, right=6mm, top=4mm, bottom=4mm,
  fontupper=\color{white}\sffamily,
  title={\sffamily\bfseries\color{white}\ding{72}~~À retenir},
  colbacktitle=qtealdark, coltitle=white, boxrule=0pt
]
\begin{itemize}[leftmargin=6mm, label={\color{qamber}\ding{212}}, itemsep=1mm]
  \item Arithmétique : $u_n=u_0+nr$, somme $=(\text{nb termes})\times\frac{\text{premier}+\text{dernier}}{2}$. Géométrique : $u_n=u_0q^n$, somme $=$ premier$\times\frac{1-q^{\text{nb}}}{1-q}$.
  \item Arithmético-géométrique $u_{n+1}=au_n+b$ : $\ell=\frac{b}{1-a}$, $v_n=u_n-\ell$ géométrique de raison $a$.
  \item Convergence monotone : croissante+majorée (ou décroissante+minorée) $\Rightarrow$ converge ; ne donne pas la valeur.
  \item Suites adjacentes : $(u_n)$ croissante, $(v_n)$ décroissante, $v_n-u_n\to0$ $\Rightarrow$ elles convergent vers une \textbf{même} limite $\ell$, avec $u_n\leq\ell\leq v_n$.
  \item Méthode alternative : $|f'|\leq k<1$ sur intervalle stable $\Rightarrow |u_n-\ell|\leq k^n|u_0-\ell|\to 0$.
  \item Rédaction d'une suite récurrente : intervalle stable (récurrence) $\to$ monotonie $\to$ convergence $\to$ résoudre $f(\ell)=\ell$.
  \item Sommes télescopiques : $\sum(a_{k+1}-a_k) = a_{n+1}-a_0$.
\end{itemize}
\end{tcolorbox}

\end{document}
